% Included by main.tex. One specification, two implementations.
\chapter{Implementations and how they compare}\label{ch:compare}

The requirements in Chapter~\ref{ch:tasks} and the specification files fix what
Maqam does; they say nothing about the language. \path{manifest.json} lists two
implementations of the same specification, each mapped to it by its own
\path{impl/<name>/Implementation.lean} and checked by \cmd{make spec-check}.

\begin{description}
  \item[\path{impl/rust/source}] the upstream Rust, the reference the
  requirements were distilled from, with one change: its reverb is the
  specification's (Section~\ref{sec:reverb}). Built with \cmd{cargo}.
  \item[\path{impl/zig/source}] the default implementation, a port to Zig 0.17 that runs on Linux and
  macOS. On Linux it needs only the Zig compiler and the ALSA development
  headers (\cmd{libasound2-dev}); on macOS only the Zig compiler and the
  system's CoreAudio, with the Command Line Tools' SDK. ffmpeg is needed at run
  time for \cmd{m} (recording), as upstream.
\end{description}

\section{Building the Zig implementation on its own}
Nothing under \path{impl/rust/} is read, built or linked by the Zig build
(proof \emph{Implementations stand alone}). The bundled \path{.mq} sessions are
copied into \path{impl/zig/source} because they are inputs of the
specification, not of the Rust build. From \path{impl/zig/source}:

\begin{center}
\begin{tabular}{ll}
\toprule
\cmd{make} & build \path{zig-out/bin/maqam-live} (ReleaseFast)\\
\cmd{make check} & build and run every module's tests\\
\cmd{make render SESSION=default.mq} & offline render to \path{render.wav}\\
\cmd{make bench SESSION=default.mq NAM=x.nam} & the real-time bench at 64 frames\\
\cmd{make run} & the terminal app\\
\bottomrule
\end{tabular}
\end{center}

\section{One device per operating system}
The engine never knows which device drives it: \texttt{device.Output} is the
ALSA device on Linux and the CoreAudio device on macOS, chosen when the program
is compiled, and each calls \texttt{audio.Engine.process} once per block. They
differ in who keeps time. ALSA is written to: the device thread renders a block
and blocks in the write until the hardware has room. CoreAudio pulls: it calls
the output callback on its own real-time thread when it needs a block, and a
second, input-only unit hands the captured frames to that callback through a
lock-free ring. The requirements on the callback are the same on both.

The controls behind \emph{The audio thread never waits on the kernel} and
\emph{A buffer skip is never silent} are the operating system's:

\begin{center}
\begin{tabular}{p{0.24\linewidth}p{0.33\linewidth}p{0.33\linewidth}}
\toprule
 & Linux (ALSA) & macOS (CoreAudio)\\
\midrule
real-time scheduling & \texttt{SCHED\_FIFO}, else RealtimeKit (\texttt{fifo}, \texttt{rtkit}) & the Mach time-constraint policy (\texttt{mach}); CoreAudio's own I/O thread already has it\\
no page faults & \texttt{mlockall}, current and future mappings & no \texttt{mlockall} on macOS (\texttt{memory\_locked} false): the callback's buffers are written once before the device starts\\
a skip is counted & the write returns \texttt{-EPIPE} & the device's sample time jumps past the frames handed over\\
opening a URL & \cmd{xdg-open} & \cmd{open}\\
\bottomrule
\end{tabular}
\end{center}

\noindent On macOS the terminal running \cmd{maqam-live} needs microphone
permission for the live input; without it the app plays with no input, as on a
Linux machine with no capture device.

\noindent From the project root \cmd{make zig-check} does the same as
\cmd{make check} there, and \cmd{make check} runs the specification check and
both implementations' checks.

\section{How the Zig is arranged}
Each Rust module has a Zig module of the same name and role; types keep their
Rust names and functions become camelCase (\path{impl/zig/source/PORTING.md}).

\begin{longtable}{lll}
\toprule
Rust & Zig & role\\
\midrule
\endhead
\texttt{tuning.rs} & \texttt{tuning.zig} & pitches, ajnas, the 5-limit lattice\\
\texttt{command.rs} & \texttt{command.zig} & the command language and its messages\\
\texttt{sequencer.rs}, \texttt{synth.rs} & same names & phrases, bars, walks, voices\\
\texttt{session\_v3.rs} + loaders in \texttt{app.rs} & \texttt{session.zig} & MAQAM\_SESSION\_V3/V2/V1\\
\texttt{fx.rs}, \texttt{vcf.rs}, \texttt{analog.rs}, \texttt{sympathetics.rs} & same names & the signal chain\\
crate \texttt{nam-rs} 0.4 & \texttt{nam.zig} & NAM WaveNet (A1, A2), LSTM, slimmable\\
\texttt{audio.rs} + \texttt{cpal} & \texttt{audio.zig}, \texttt{device.zig}, \texttt{coreaudio.zig} & the engine, ALSA and CoreAudio\\
\texttt{record\_old.rs}, \texttt{carpet.rs} & \texttt{record.zig}, \texttt{carpet.zig}, \texttt{png.zig} & the MP4 recorder\\
\texttt{app.rs} & \texttt{app.zig}, \texttt{net.zig} & the App and its workers\\
\texttt{ui.rs} + \texttt{ratatui} & \texttt{ui.zig}, \texttt{term.zig} & the terminal UI\\
crossbeam channels, main.rs statics & \texttt{channel.zig}, \texttt{shared.zig} & lock-free rings, atomics\\
--- & \texttt{eventlog.zig}, \texttt{wav.zig} & the comparison hooks\\
\bottomrule
\end{longtable}

The audio thread in the Zig never allocates, frees or locks. The App sends
pointers to phrases and models on a 512-slot lock-free ring; the engine sends
what it no longer uses back on a second ring and the App frees it. Upstream's
mutex-guarded jump-counter map is a fixed table of atomics.

\section{Where the Zig differs on purpose}
Messages, ranges, session files and commands are the same. The differences are
the requirements the Rust does not meet (Chapter~\ref{ch:gaps}):

\begin{itemize}
  \item A phrase whose combined scale has fewer than three notes is rejected
  with a message instead of panicking (\cmd{create Zz 1/1}, \cmd{d zz 4}).
  \item Inputs that loop forever or crash upstream are handled: a non-ASCII
  value token, deeply nested \cmd{edit}/\cmd{insert}, an infinite or zero
  frequency in the tuning lattice or the sympathetic strings, and NAM files
  that make \texttt{nam-rs} index out of bounds.
  \item Jins name ties resolve alphabetically instead of by hash order, so
  \cmd{d bay} always means the same jins.
  \item Bounds: repeats at most \texttt{r20}, jump times at most 10000, the
  timeline at most 4096 rows, LLM responses at most 4\,MiB, NAM downloads at
  most 256\,MiB, NAM files at most 64\,MiB, LLM context at most 16
  exchanges, line history at most 1000 lines, recordings at most 600\,s, and
  the voice lists, jump counters and garbage queues fixed in advance.
  \item A recording that fails leaves no \path{.partial.mp4} and no temporary
  layers; after an ffmpeg failure only the \path{maqam-ffmpeg.log} the message
  points to is kept.
  \item Every NAM cache name is sanitized; tests never reach the network or
  open a browser.
  \item MIDI clock in and out (\cmd{clockin}, \cmd{clockout}), which the
  specification does not mention, answer that MIDI is not available.
  \item The audio thread locks its memory (Linux) and asks for real-time
  scheduling before its loop, and the live device counts every underrun and shows it in
  the header; upstream recovers from an underrun silently.
  \item One regression: Zig 0.17's HTTP client has no request timeouts, where
  upstream gives up after 30\,s (LLM) and 120\,s (downloads). Only the worker
  thread waits; the audio and the terminal stay live.
  \item The comparison hooks (\cmd{--render}, \cmd{--bench},
  \texttt{MAQAM\_EVENT\_LOG}) exist.
\end{itemize}

\section{Measuring them the same way}
\cmd{make compare} in the project root runs each implementation's
\cmd{check}, \cmd{render} (twice, to test determinism) and \cmd{bench}, checks
the event log against protocol \texttt{MaqamLive} with LeanFM's
\texttt{session\_conform}, checks every block of the bench's \texttt{AudioPeriod}
stream (exactly, and its conformance; Section~\ref{sec:deadline}), null-tests each render against the default
implementation's, and writes \path{reports/compare/}. A target an
implementation does not provide is reported as \emph{not provided}, never as a
pass; upstream Rust provides no hooks, so its render and bench columns say so.

\IfFileExists{../reports/compare/summary.tex}{% Generated by scripts/compare.py; do not edit.
\noindent Measured by \cmd{make compare}: session \path{impl/zig/source/default.mq}, 2 cycles at 48000\,Hz; bench 30\,s of 64-frame blocks at 48\,kHz with the sympathetic strings and the reverb, NAM \path{impl/zig/source/testdata/nam/marshall_jcm800.nam}.

\begin{longtable}{lp{0.78\linewidth}}
\toprule
\multicolumn{2}{l}{\textbf{zig (default)}}\\
specification & validated\\
check & pass 0.1 s\\
render & pass, 14.8 s, deterministic\\
event log & conforms\\
bench & 22500 blocks, 0 deadline misses; p50 218 us, p99 301 us, max 791 us; slack 542 us; scheduling rtkit, memory locked\\
every block & 0 of 22500 late, 0 labels disagree; AudioPeriod: conforms\\
source lines & 27572\\
\midrule
\bottomrule
\end{longtable}
}{%
\noindent Run \cmd{make compare} to fill this section with measurements.}
